\begin{center}
    {\large \textbf{Giorno 14} \textbar\ Oggi organizza Elisa! \textbar\ 22 Agosto 2024 \textbar\ \textbf{Ubud}, Bali}
\end{center}

\noindent\rule{\textwidth}{0.4pt}
\begin{figure}[h!] % Portrait images below
    \centering
    % First portrait image (on the left)
    \begin{minipage}[h]{0.48\textwidth} % Set width equal to half the page
        \centering
        \includegraphics[width=\textwidth]{images/flags/Screenshot sukawati.png}
    \end{minipage}
    \hfill
    % Text
    \begin{minipage}[h]{0.48\textwidth}
        \raggedright
        \textbf{Superficie}: 4.567 km² \\
        \vspace{0.3cm}
        \textbf{Popolazione}: 3,3 milioni \\
        \vspace{0.3cm}
        \textit{L'isola}
    \end{minipage}
\end{figure}

\landscapeLargeMargin{images/day14/PXL_20240822_002720316.jpg}

\landscapeLargeMargin{images/day14/PXL_20240822_002919329.jpg}

\landscapeLargeMargin{images/day14/PXL_20240822_024009362.MP.jpg}

\landscapeLargeMargin{images/day14/PXL_20240822_035948963.MP.jpg}

\landscapeLargeMargin{images/day14/PXL_20240822_040501136.MP.jpg}

\landscapeLargeMargin{images/day14/PXL_20240822_042736700.jpg}

\landscapeLargeMargin{images/day14/PXL_20240822_075046931.jpg}

\landscapeLargeMargin{images/day14/PXL_20240822_080105032.jpg}

\landscapeLargeMargin{images/day14/PXL_20240822_080637057.jpg}

\landscapeLargeMargin{images/day14/PXL_20240822_112116874.NIGHT.jpg}

\landscapeLargeMargin{images/day14/PXL_20240822_131140567.jpg}

\noindent 22 Agosto 2024 \newline

Stamattina abbiamo posticipato la colazione alle 8:30 vista la giornata pesante di ieri. Alle 8:20 la ragazza del servizio in camera è già a bussare alla nostra porta, pronta con il vassoio in mano. Elisa è ancora sotto la doccia e io faccio appena in tempo a indossare i pantaloni e ad aprirle ancora in pigiama, con gli occhi accecati dal sole — sgamato in pieno — per dirle che tempo due minuti saremmo arrivati.

Manteniamo la promessa e arriviamo incredibilmente puntuali. Quando usciamo dalla stanza vediamo in lontananza la ragazza che ha apparecchiato il tavolino con la nostra colazione e ci aspetta a fianco, in ginocchio, pronta a spiegarci il menù. Ormai mi vergogno come un ladro di averle aperto in pigiama, ma il danno è fatto.

Questa volta abbiamo cercato di non commettere errori: caffè e frullato di frutta per entrambi, mentre come piatto Elisa ha ordinato un Avocado Smash — due bruschette con avocado, formaggio, pesto, cardamomo, pomodoro, aglio, cipolla, un filo d'olio e cetriolo — e io un Akuri Tofu with Creamy Slaw Breakfast Burger, un hamburger di tofu per dirla semplice, che a me sembra un panino con tutti i condimenti in cui hanno dimenticato di mettere la carne. Il commento finale sulla colazione è: "Ma che ci vuole a fare due uova?" Si tratta però della nostra ultima colazione al Gayatri e siamo un po' malinconici, quindi tutto ci va bene.

Finito di mangiare torniamo in camera a preparare le borse per il check-out, ma vogliamo goderci quel paradiso fino all'ultimo: una volta preparato tutto andiamo a sdraiarci sui lettini a bordo piscina. Eli fa anche il bagno con tanto di servizio fotografico ad opera del sottoscritto, mentre io cerco di rilassarmi sul lettino, rigorosamente all'ombra.

Che siamo giunti a fine viaggio lo si capisce anche dal fatto che Elisa inizia il suo rito ormai ben collaudato di fare la lista delle cose che non abbiamo visto — una roba che, per chi ha organizzato tutto, sembra tanto una pugnalata al cuore. Oggi il piano lo fa lei: come prima tappa a fine mattinata Sukawati, un villaggio a sud di Ubud famoso per i mercatini locali con souvenir caratteristici a prezzi stracciati, e nel pomeriggio il Tempio Madre di Besakih, che a detta di molti è uno dei migliori di tutta Bali.

Terminata la pacchia in piscina facciamo il check-out, approfittando ancora del posto e lasciando i bagagli fino a sera. Rifiutiamo però l'aiuto del manager della struttura — un giovane alto con lineamenti misti orientali-europei che potrebbe fare il modello — che si offre generosamente di portarmi lo zaino fino alla reception. Lo abbiamo visto girare per la struttura più volte, controllando e impartendo, educatamente ma con fermezza, ordini ai dipendenti per rendere tutto efficientissimo. Anche adesso non manca di essere d'aiuto, ma direi che dieci metri con lo zaino in spalla posso farli anche da solo.

Lasciamo il Gayatri e un Grab ci sta già aspettando per portarci a Sukawati.

\begin{center}
    % Decorative line
    \noindent\rule{0.6\textwidth}{0.4pt}
\end{center}

Quello che troviamo all'arrivo non è esattamente ciò che ci aspettavamo. Ci immaginavamo mercati all'aperto come i souk israeliani, con bancarelle di artigianato, arte e abbigliamento dappertutto. Percorriamo invece una via principale che, anche a detta dell'autista del Grab, dovrebbe essere quella del mercato, e troviamo solo venditori di cibo — chiamarli bar mi sembra eccessivo — e qualche bancarella di souvenir molto turistici. Guardando su Maps troviamo una zona adibita a mercato ma quando la esploriamo non troviamo nulla, anche se almeno ci permette di fare un giro dietro le case e spiare un po' la vita locale.

Arriviamo poi a un crocevia dove si trova un complesso di palazzi di tre piani che prende il nome di Sukawati Art Market. Uno dopo l'altro visitiamo gli edifici, che sono in sostanza dei centri commerciali con luci al neon e telecamere a ogni piano, allestiti all'interno con bancarelle turistiche. A ogni bancarella siamo assaliti dalle venditrici che ci propongono la stessa merce scadente che si trova dappertutto: magliette della birra Bintang stampate male, Sarong sintetici e cazzi di legno. Questa cosa dei cazzi di legno sta decisamente sfuggendo di mano — ormai ogni paese del mondo li propone come souvenir insieme ai prodotti tipici, non capisco perché mai qualcuno dovrebbe tornare da una vacanza portando come ricordo un cazzo di legno. A parte forse i francesi.

La pressione delle venditrici ci fa sempre più venir voglia di andarcene e così è. Terminati i palazzi troviamo un mercato all'aperto che sembra avere prodotti più tipici, ma è solo un'illusione. Riesco a provare un paio di pantaloni — che scopro poi essere da donna, che non mi piacciono nemmeno — per i quali sparano un prezzo altissimo per poi abbassarlo prima ancora che inizi la contrattazione.

Ce ne andiamo scoraggiati e percorriamo invece un mercato alimentare locale, dove per fortuna non cercano di farci abboccare. Qui siamo colpiti da un venditore di pulcini che ne possiede un esercito in gabbiette. Tra di loro vi sono pulcini colorati chimicamente — verdi, rossi e giallo fluo — che sembrano aver visto giorni migliori. Proviamo a chiedergli come abbia fatto a ridurli in quello stato, ma non parla una parola di inglese. Scopriremo poi che è una pratica sempre più in uso: i pulcini vengono colorati iniettando coloranti nelle uova o spruzzandoli direttamente una volta nati.

Quella è la goccia che fa traboccare il vaso. Dichiariamo l'esperienza a Sukawati conclusa e chiamiamo subito un Grab per il tempio di Besakih.

\begin{center}
    % Decorative line
    \noindent\rule{0.6\textwidth}{0.4pt}
\end{center}

Durante il tragitto Eli si addormenta e stavolta sono costretto a svegliarla: russa. Il tempio si trova lontano dalle città, raggiungibile dopo vari tornanti in una strada che passa per la foresta. Situato in cima a una salita, dispone di un ampio parcheggio multipiano e di un'area di accettazione che sembra quella di un aeroporto. È anche il più caro di tutti, ma insieme al biglietto abbiamo a disposizione una navetta che ci accompagna fino al tempio, una guida per il tour e ovviamente il solito Sarong.

Vediamo però subito come tutto questo pacchetto sia un mezzo scam per i turisti. Il Sarong è in prestito ed è il più brutto di tutti quelli provati, un telo sintetico grigio tinta unita. La navetta ci accompagna in salita ma non in discesa — fatto apposta perché al ritorno si cammini tra le bancarelle di souvenir. Se avessero messo il banco dell'accettazione all'effettivo ingresso del tempio la navetta non sarebbe stata necessaria. In ultimo la guida: il tour dura una mezz'ora, compresi dieci minuti in coda per la foto di rito, e a noi stavolta va particolarmente male.

Arrivati in cima scopriamo di essere con altri italiani e di avere due guide, una in inglese e una in italiano. Mi viene subito il dubbio che il gruppo di italiani abbia una guida privata e che quella assegnata a noi sia quella in inglese, ma Eli è contenta di avere per una volta una guida che parla italiano e così decido di non farmi troppe domande e seguiamo il gruppo. La guida inglese nel frattempo si dilegua.

Andando avanti scopriamo che avevo visto giusto: gli italiani hanno acquistato un tour di tre giorni con guida privata. Non sembrano avere problemi al fatto che seguiamo con loro, salvo che dopo due chiacchiere di rito ci congedano con un "allora ciao e buon proseguimento." Va bene, ce ne andiamo, godetevi la vacanza.

Mi hanno raccontato il loro itinerario: tre giorni di templi, poi cinque giorni sulla Gili Trawangan e poi altri tre in una città di Bali. Penso che neanche impegnandomi avrei potuto organizzare una cagata simile: tre giorni di corsa per i templi con una guida incapace, cinque giorni sulla Party Gili che per un gruppo con una bambina piccola è la scelta peggiore delle tre isole, e il resto in una città commerciale di Bali — probabilmente in qualche centro balneare dove tra due anni guarderanno le foto chiedendosi: "Ma qui eravamo in Indonesia o a Spotorno?" Forse sono un po' cinico perché ci siamo presi la fregatura della guida pagata ma non usufruita, però era evidente che sapesse poco del tempio: ha passato i primi dieci minuti a tradurre un'insegna col nome del tempio spiegandoci come funzionasse la scrittura balinese.

Proviamo allora a dare qualche informazione realmente utile almeno qui.

\begin{center}
    % Decorative line
    \noindent\rule{0.6\textwidth}{0.4pt}
\end{center}

\textit{Il tempio di Besakih è un complesso situato nel villaggio di Besakih, distretto di Rendang, reggenza di Karangasem. È composto da un tempio centrale — il Penataran Agung Besakih — e 18 templi complementari, tra cui il Tempio Basukian. Quest'area fu il primo luogo in cui la rivelazione di Dio fu ricevuta da Hyang Rsi Markendya, il precursore della religione indù del Dharma a Bali. Il Penataran Agung è il tempio più grande del complesso, con il maggior numero di edifici e cerimonie, e ospita tre santuari principali chiamati Padma, dedicati ai tre livelli di coscienza spirituale del Tri Purusha: Shiva, Sada Shiva e Parama Shiva.}

\begin{center}
    % Decorative line
    \noindent\rule{0.6\textwidth}{0.4pt}
\end{center}

La frase giusta qui sarebbe: "Bello, ma il mio falegname con 30.000 rupie lo faceva meglio." Non ci rimane particolarmente impresso, né come architettura né come esperienza spirituale. Il motivo vero però è probabilmente lo stesso che ci aveva fatto rimuovere il tempio di Lempuyang dall'itinerario: la foto con lo specchio.

Lempuyang è diventato forse il tempio più famoso di Bali grazie a una foto che spopola su Instagram: la coppietta in posa al centro della grande porta, con una montagna sullo sfondo che si riflette ai loro piedi in uno specchio d'acqua. Salvo che ai loro piedi non c'è acqua: l'effetto è ottenuto dal fotografo che pone uno specchietto alla base della fotocamera. Questo ha fatto sì che tutti volessero andare a Lempuyang per ritrarsi "riflessi", con code normalmente fino a tre ore. Qui a Besakih la situazione è analoga: alla base del tempio centrale c'è la fila per farsi la foto e i locali chiedono 30.000 rupie — casualmente le stesse del mio falegname — per scattartela con lo specchio. Non ho problemi a dire, quando ci propongono il pacchetto fotografico, che "we don't want fake photos", e così ci facciamo scattare delle foto normali dalla coppia di turisti dietro di noi.

Procediamo la visita cercando di captare le informazioni di qualche guida, ma di informazioni ce ne sono davvero poche. Di guide, invece, è pieno. Usciti dal tempio scendiamo lo stradone lungo il quale eravamo saliti con la navetta e non troviamo souvenir che ci attraggano, tutto troppo turistico. Non siamo ancora riusciti a pranzare ma quando chiediamo il prezzo di un pacco di patatine in una bancarella per poco non gli scoppiamo a ridere in faccia.

\begin{center}
    % Decorative line
    \noindent\rule{0.6\textwidth}{0.4pt}
\end{center}

A quel punto inizia l'odissea. Trovare un Grab o un Gojek in un posto sperduto come questo non è facile: un conto è trovarli in città per andare fuori, un altro è il contrario. Quando un driver fa il match con la mia chiamata sono sorpreso, salvo scoprire dalla foto sull'app che si tratta di un ragazzino — sicuramente maggiorenne, ma con la faccia da lattante — che dopo aver confermato la corsa mi scrive che il prezzo dell'app è basso per il percorso che gli sto chiedendo. Gli rispondo seccato: "I don't make the price." Mi chiede se sia disposto ad aggiungere qualcosa e la mia risposta è un immediato no. Cancello la corsa.

Riprovo con Gojek ma la situazione è la stessa: nessun driver nella zona, se non il ragazzino che dopo aver capito che sono sempre io mi cancella a sua volta la corsa. Continuo a provare per venti minuti buoni, seduto in mezzo a un incrocio, allontanando di volta in volta i motorini che si avvicinano chiedendomi se voglia un taxi. Intanto il sole si abbassa e la fame sale.

Decidiamo di tornare al parcheggio del tempio e chiediamo agli addetti alla sicurezza se tramite loro sia possibile chiamare un taxi. Sapendo che un taxi da lì fino a Ubud mi avrebbe chiesto una cifra spropositata, gli dico che mi serve arrivare a un centro abitato a metà strada — Banjarangkan, scelto a caso sulla mappa — da cui ho intenzione di prendere un Grab. A loro spiego invece che lì mi aspettano degli amici pronti a darmi un passaggio.

I tre della sicurezza si consultano velocemente e il più giovane si propone di accompagnarci per la modica cifra di 30.000 rupie — numero che in questa giornata ritorna un po' troppo spesso. Sono incredulo — si tratta di circa 2 euro — e chiedo conferma facendo anche il segno del tre con le dita, perché la cifra è bassissima rispetto a quanto proposto dalle app. Accettiamo subito.

Durante il tragitto non parliamo. Elisa dorme mentre io, col fatto di non essere su un mezzo ufficiale verificato come quelli di Grab o Gojek, rimango tutto il tempo a monitorare il percorso su Maps. Ovviamente è tutta una paranoia mia e lui procede dritto alla meta. Provo anche il colpaccio: gli dico che i nostri amici non ci possono più dare il passaggio e gli offro di pagarlo di più se ci avesse portati fino a Ubud. Mi spiega, con molta educazione ma in un inglese stentato, che dopo deve tornare a lavorare e quindi non può.

Appena arriviamo alla cittadina ci fermiamo di fronte a un supermercato e gli chiedo se ha il resto da darmi, mostrandogli una banconota da 100.000 rupie. Mi guarda stranito e a stento mi fa capire che di quelle banconote ne vuole tre. Rimango di stucco: gli dico che gli accordi non erano quelli, che non ho intenzione di farmi fregare e che gli avevo chiesto esplicitamente per ben due volte se la cifra fosse effettivamente 30.000 rupie.

Lui tira fuori la calcolatrice e scrive 300 continuando a dire "thirty", sempre con tono educato. Inizio a pensare che non stesse cercando di fregarmi ma che fosse davvero un disastro con l'inglese — e infatti passa dalla calcolatrice al traduttore. Dal traduttore leggo che si scusa, che sta studiando l'inglese e che non lo parla molto bene, ma che la cifra chiesta erano 300.000 rupie e non 30.000.

Tramite il traduttore indonesiano scambiamo ancora un paio di messaggi senza arrivare a nulla, finché decido di contrattare. Gli scrivo che ho capito che non sa l'inglese, ma che io ho accettato una cifra e lui ora me ne sta dicendo un'altra e la colpa non è mia. Vorrei proporre 150.000, che sarebbe la cifra giusta per la strada percorsa, ma avendo solo banconote da cento mi spingo fino a 200.000 per chiuderla, e lui accetta. Quando scendiamo dall'auto si scusa ripetendo a stento che sta ancora studiando l'inglese. "Keep studying" è la mia risposta uscendo.

Sicuramente l'inglese non lo parlava bene, ma come scam ai turisti ci sapeva fare — se non lui, che effettivamente sembrava onesto, almeno i suoi colleghi della sicurezza, davanti ai quali avevo ripetuto la domanda sul prezzo appena me lo aveva proposto.

Dalla stazione di servizio troviamo subito un Grab disposto a portarci a Ubud e mentre lo aspettiamo compriamo patatine e gelato per riempire il vuoto nello stomaco. Anche il nostro autista pare aver avuto un'idea simile poco prima, perché non appena saliamo scende per buttare gli incarti di due panini tipo McDonald's. È giovane pure lui e sembra il Piccolo Lucio indonesiano, con tanto di faccione con gli occhiali spessi.

Facendo due conti, tra la cifra data all'improvvisato tassista e quella data al Grab, abbiamo speso più o meno come se avessimo trovato un passaggio direttamente dal tempio. Quindi non mi sento di dire di essere stato truffato, e soprattutto sono contento di non aver ceduto al ricatto del ragazzino che mi aveva negato il passaggio sull'app.

\begin{center}
    % Decorative line
    \noindent\rule{0.6\textwidth}{0.4pt}
\end{center}

Approfitto di questo ultimo tratto in macchina per organizzare il viaggio dell'indomani. Mi sono già occupato dei check-in aerei per i voli di sabato e adesso ricontatto il nostro caro Ketut chiedendogli se può portarci a Kuta la mattina seguente. Provo a contrattare un po' ma capisco che questi tragitti tra città hanno prezzi standard e chiudiamo presto la trattativa.

Ormai si è fatto buio e quando arriviamo a Ubud è ora di cena. Entriamo per l'ultima volta al Gayatri a ritirare i bagagli, scattiamo le ennesime fotografie del posto e lo salutiamo per sempre, spostandoci nell'albergo a fianco prenotato all'ultimo per tappare l'ultima notte a Ubud. Il costo della nuova sistemazione è circa un decimo rispetto al Gayatri. Li vale tutti.

Entriamo in quello che sembra un cantiere, con tanto di montagne di sabbia. Attraversiamo un sentiero dissestato che passa di fianco a una piscina di pietra che ricorda quella in Messico "venduta già col cadavere dentro" e arriviamo a una reception improvvisata di fronte a una stanzina, da cui esce un ragazzo scalzo che ci consegna la chiave della stanza.

La stanza ha tutte e quattro le mura, questo si può dire, ed è molto spaziosa. Non ci sono però altre cose positive da raccontare: il bagno puzza, lo scarico è rotto, non c'è l'aria condizionata ma solo una ventola a soffitto che fa un rumore allucinante, le pareti sono sottilissime e si sente tutto ciò che succede fuori. Amen — dobbiamo passarci una notte e vista la posizione era la scelta più strategica, dal momento che il Gayatri non aveva più disponibilità.

Ci docciamo, lasciamo i bagagli e usciamo, andando finalmente a provare il ristorante di fronte al quale passavamo ogni sera ripromettendoci di andarci. Il locale è isolato da Monkey Street, all'aperto, con musica live. Il gruppo che suona è particolarmente intonato e spazia su più generi musicali, con una predisposizione evidente per i Coldplay.

Come ordinazione andiamo sul collaudato: Mie Goreng con involtini primavera, Bintang e Bintang Radler. Finita la cena cambiamo tavolo e ordiniamo un'altra birra per assistere più da vicino alla live. Sono stanchissimo — ho fatto meno pisolini in auto oggi — e vogliamo ancora comprare qualche souvenir, così dopo qualche canzone decidiamo di alzarci. Alla cassa sono molto distratti e per poco non mi lasciano uscire avendo pagato solo l'ultima birra, ma facciamo gli onesti e li avvisiamo.

Per i souvenir andiamo niente meno che al Coco Market, per la terza sera da quando siamo a Ubud. Prendiamo caffè, tè, spezie, zuppe, zucchero di cocco e altri prodotti alimentari, un po' per noi e un po' da regalare. Terminati gli acquisti torniamo di buona lena in camera e alla fine avremmo probabilmente bisogno di un'altra doccia, ma la stanchezza è troppa e l'indomani ci aspetta un altro viaggio. Andiamo subito a dormire.

\pagestyle{empty} % disable numbers

%coppia
\landscapeThinMargin{images/day14/PXL_20240822_002433561.jpg}
\landscapeThinMarginCentered{images/day14/PXL_20240822_024239071.MP.jpg}
%

%coppia
\landscapeThinMarginCentered{images/day14/PXL_20240822_024023334.MP.jpg}
\landscapeThinMarginCentered{images/day14/PXL_20240822_024019214.MP.jpg}
%

%coppia
\landscapeThinMarginCentered{images/day14/PXL_20240822_024115158.MP.jpg}
\landscapeThinMarginCentered{images/day14/PXL_20240822_024110551.jpg}
%

%coppia
\landscapeThinMarginCentered{images/day14/PXL_20240822_041908561.jpg}
\landscapeThinMarginCentered{images/day14/PXL_20240822_072107180.jpg}
%

%coppia
\splitImageLeft{images/day14/PXL_20240822_075114469.MP.jpg}
\splitImageRight{images/day14/PXL_20240822_075114469.MP.jpg}
%

%
\portraitFullscreen{images/day14/PXL_20240822_124548341.NIGHT.jpg}
%

\pagestyle{fancy} % enable numbers